\documentclass[12pt]{article}
\usepackage[utf8]{inputenc}
\usepackage[T1]{fontenc}
\usepackage[brazil]{babel}
\usepackage[a4paper,margin=2.5cm]{geometry}
\usepackage{enumitem}
\setlist[itemize]{noitemsep,topsep=4pt}
\title{Sprint 3 - S3\_14\\Portas Lógicas e Arquitetura}
\author{Matheus Moreli Durão Silva\\Matrícula 12621CBS217\\Cibersegurança - Bacharelado - Noturno}
\date{2026}
\begin{document}
\maketitle

\section*{Entregável}
Relatório de álgebra booleana e arquitetura de Von Neumann. Este documento registra o objetivo do card, a solução desenvolvida e a forma de conferência. Os arquivos de código e as evidências citadas permanecem na mesma pasta deste documento.

\section*{O que foi estudado e aplicado}
\begin{itemize}
\item Portas AND, OR e NOT
\item Tabelas-verdade
\item Circuitos combinacionais
\item CPU, memória e barramentos
\end{itemize}

\section*{Como conferir}
O avaliador pode abrir os arquivos desta pasta, executar os exemplos quando houver código e comparar a saída com a explicação apresentada. Nos cards ligados ao Inventário de Ativos, a solução foi relacionada ao Trabalho 2 sem modificar requisitos já atendidos do CRUD.

\section*{Portas e circuito combinacional}
O script \texttt{portas\_logicas.py} gera tabelas para AND, OR, NAND, NOR, XOR e XNOR. Um circuito de acesso usa $S=(C\land P)\land\neg B$, com credencial $C$, permissão $P$ e bloqueio $B$: uma AND combina C/P, uma NOT inverte B e a AND final decide. Das oito combinações, apenas $(1,1,0)$ libera. A absorção $A\land(A\lor B)=A$ é verificada para todas as entradas, eliminando duas operações no trecho simplificado.
\section*{Arquitetura de Von Neumann}
A memória guarda dados e instruções; a unidade de controle busca e decodifica; a ULA executa operações; registradores guardam valores imediatos; barramentos transportam dados, endereços e sinais. No ciclo busca--decodificação--execução, o contador de programa aponta a próxima instrução, que a CPU lê, interpreta e executa. Condições em software são implementadas por comparações e desvios apoiados em circuitos lógicos. A análise completa está em \texttt{resolucao\_portas.md}.
\section*{Conclusão}
O exercício conecta álgebra booleana, circuito combinacional e execução de programas.
\end{document}
